\documentclass[
    12pt,
    a4paper,
    twoside,
    openright, % Kapitel beginnen rechts
    % draft
]{report}

% Schriftzeichen Encoder
\usepackage[utf8]{inputenc}

% Deutsch
\usepackage[ngerman]{babel}

% English
%\usepackage[english]{babel}

% Eigene Befehle
\input{settings/commands}

% Referenzen
\usepackage{cite}

% Mathematik
\usepackage{amsmath}
\usepackage{amsfonts}
\usepackage{amssymb}
\usepackage{bm}
\usepackage{bbm}        % Blackboard Schrift (Reelle Zahlen \mathbb{R})
\usepackage{amsthm}     % Mathematische Umgebungen (Definitionen, Theoreme, ...)
\usepackage{wasysym}    % Durchmesser Zeichen (\diameter)

% Grafiken
\usepackage{graphicx}
\usepackage{epstopdf}
\usepackage{tikz}
\usepackage{pgfplots} % erstellen von Plots mit Latex
\pgfplotsset{compat=1.12}
\usetikzlibrary{external}
\tikzexternalize[prefix=plots/]

% Seitenlayout
\usepackage{fancyhdr}
\setlength{\headheight}{6mm}
\setlength{\textheight}{210mm}
\setlength{\textwidth}{140mm}
\setlength{\oddsidemargin}{9.6mm}
\setlength{\evensidemargin}{9.6mm}
\setlength{\footskip}{10mm}
\setlength{\paperwidth}{210mm}
\setlength{\paperheight}{297mm}

\pagestyle{fancy}
\renewcommand{\chaptermark}[1]{\markboth{\chaptername\ \thechapter.\ #1}{}}
\renewcommand{\sectionmark}[1]{\markright{\thesection.\ #1}}
\fancyhead[LE,RO]{\thepage}
\fancyhead[RE]{\leftmark}
\fancyhead[LO]{\rightmark}
\fancyfoot[C]{}

% Tabellen
\usepackage{array}
\usepackage{booktabs}
\usepackage{rotating}

% Entfernt Kopf- und Fußzeilen auf leeren Seiten
\usepackage{emptypage}

% Zeilenabstand in itemize Umgebung
\newlength{\itemizeskip}
\setlength{\itemizeskip}{0.2em}

% Aufzählungszeichen
\def\labelitemi{-}
\def\labelitemii{$\diamond$}
\def\labelitemiii{\textbullet}

% Einheiten
\usepackage[binary-units=true]{siunitx}

% Dynamische Verknüpfungen
% ! Dieses Packet muss als letztes eingebunden werden
\usepackage[pdfborder={0 0 0}]{hyperref} 